% Generated draft: requires mathematical and source review.
\documentclass[11pt]{article}
\usepackage[margin=25mm]{geometry}
\usepackage{fontspec}
\setmainfont{Latin Modern Roman}
\usepackage{amsmath,amssymb,mathtools,unicode-math}
\setmathfont{Latin Modern Math}
\usepackage{enumitem,needspace,xurl,hyperref}
\setlength{\parindent}{0pt}
\setlength{\parskip}{5pt}
\setlength{\emergencystretch}{3em}
\newcommand{\N}{\mathbb{N}}
\newcommand{\Z}{\mathbb{Z}}
\newcommand{\Q}{\mathbb{Q}}
\newcommand{\R}{\mathbb{R}}
\newcommand{\C}{\mathbb{C}}
\newcommand{\F}{\mathbb{F}}
\newcommand{\A}{\mathbb{A}}
\newcommand{\PP}{\mathbb{P}}
\newcommand{\E}{\mathbb{E}}
\newcommand{\eps}{\varepsilon}
\newcommand{\dd}{\,\mathrm{d}}
\begin{document}
\section*{UnsolvedMath ID: 700150}
\textbf{Source problem code:} catalogue entry
\subsection{Definitions and mathematical statement}
For a group \(G\) with identity \(1_G\) and a subgroup \(H\leq G\), define its normalizer by
\[
N_G(H)=\{g\in G:gHg^{-1}=H\}.
\]
Using the standard interpretation of the source terminology, \(G\) satisfies the \emph{normalizer condition} if
\[
H\subsetneq N_G(H)\qquad\text{for every subgroup }H\subsetneq G.
\]
The supplied record does not explicitly define this condition; this interpretation requires confirmation against the source.

Is it true that every group \(G\neq\{1_G\}\) satisfying the normalizer condition has a subgroup \(A\leq G\) such that
\[
A\neq\{1_G\},\qquad
 gAg^{-1}=A\quad\text{for every }g\in G,\qquad
 ab=ba\quad\text{for all }a,b\in A?
\]
No finiteness assumption on \(G\) is imposed.
\subsection{Short English statement}
Must every nontrivial group in which each proper subgroup is strictly contained in its normalizer contain a nontrivial abelian normal subgroup? This uses the standard interpretation of the source's undefined phrase ``normalizer condition.''
\subsection{Sources}
\begin{itemize}
\item Local UnsolvedMath registry, supplied complete source record: permanent ID \(700150\), exact source problem code \texttt{catalogue entry}, title ``Kourovka Notebook Problem 2.80.'' The registry labels the problem open; this label is supplied metadata, not an independently verified assessment.
\item \emph{The Kourovka Notebook (Unsolved Problems in Group Theory)}, Problem 2.80, proposed by S. N. Chernikov, 2nd issue (1967). The supplied record identifies the 21st issue edition (2026), \texttt{arXiv:1401.0300v48}. The supplied collection description names E. I. Khukhro and V. D. Mazurov as editors. These bibliographic details are supplied attribution; the cited edition has not been independently consulted.
\end{itemize}
\end{document}
